% Physical unit 006; accepted Session 02.
\GordanUnitStart{006}{printed page 1}
\begin{center}
{\large CHAPTER I.}\\[0.8em]
\S\ 1.
\end{center}

A spheroid is a surface formed by the rotation of an ellipse about its minor axis \gq{$b$}. In this motion the individual points of the periphery of the ellipse describe circles (parallels), the greatest of which, described by the extreme points of the major axis \gq{$a$}, is called the equator. ---

\begin{center}\S\ 2.\end{center}
The minor axis of the ellipse, about which the rotation has been made, is called the axis of the spheroid; its extreme points are the poles; and every plane containing that axis is a meridian.

\begin{center}\S\ 3.\end{center}
The angle \gq{$\psi$}, namely the one formed by the meridian containing the point \gq{$A$} and any chosen first meridian, is called the longitude; but the angle \gq{$B$}, which the normal erected at $A$ on the spheroid makes with the plane of the equator, is the latitude of the point $A$. Hence one sees that every point of the spheroid is determined by longitude and latitude. ---

\begin{center}\S\ 4.\end{center}
Since all points of the same longitude $\psi$ lie on the same meridian, and all points of the same latitude $B$ lie on the same circle, we call the angle $\psi$ the longitude of the meridian containing those same points, and the angle $B$ the latitude of the circle containing those same points. ---

\begin{center}\S\ 5.\end{center}
Besides the latitude $B$, let us use another angle, \gq{$l$}, to determine the position of a circle. For, if in place of the radius of the spheroid we use the axis \gq{$a$} and describe a sphere about the center, the plane of a circle of the spheroid cuts this sphere also in a circle. Each of the normal lines erected on the sphere at the points of the periphery
\GordanBottomFolio{1}
